\documentclass[10pt,a4paper]{article}
\usepackage[margin=12mm]{geometry}
\usepackage{fontspec}
\usepackage{xeCJK}
\setmainfont{texgyretermes-regular.otf}[BoldFont=texgyretermes-bold.otf,ItalicFont=texgyretermes-italic.otf,BoldItalicFont=texgyretermes-bolditalic.otf]
\setCJKmainfont{FandolSong-Regular.otf}[BoldFont=FandolSong-Bold.otf]
\usepackage{enumitem,tabularx,titlesec,xcolor}
\usepackage[colorlinks=true,urlcolor=blue!45!black,linkcolor=blue!45!black]{hyperref}
\hypersetup{pdftitle={陈俊乐 - 个人简历},pdfauthor={Junle Chen},pdfsubject={Academic Curriculum Vitae}}
\pagestyle{empty}
\setlength{\parindent}{0pt}
\setlength{\parskip}{0pt}
\setlength{\tabcolsep}{4pt}
\setlength{\emergencystretch}{2em}
\titleformat{\section}{\fontsize{11.5}{15}\selectfont\bfseries}{}{0pt}{}[\vspace{1.5pt}\titlerule]
\titlespacing*{\section}{0pt}{6pt}{4pt}
\begin{document}
\fontsize{9.7}{11.2}\selectfont
\begin{center}
{\fontsize{24}{32}\selectfont\bfseries 陈俊乐\par}\vspace{4pt}
{\fontsize{12}{18}\selectfont 计算机科学与工程博士生\par}\vspace{3pt}
香港\enspace|\enspace \href{mailto:jchenkg@connect.ust.hk}{jchenkg@connect.ust.hk}\enspace|\enspace \href{tel:+8615259695906}{+86 15259695906}\par\vspace{2pt}
\href{https://junle-chen.github.io}{junle-chen.github.io}\enspace|\enspace \href{https://github.com/junle-chen}{GitHub: junle-chen}\enspace|\enspace \href{https://scholar.google.com/citations?user=EjiEmLUAAAAJ&hl=zh-CN}{Google Scholar}\enspace|\enspace 微信：15259695906
\end{center}
\vspace{-8pt}
\section*{研究概况}

香港科技大学计算机科学与工程系博士生，导师为周晓方教授。研究聚焦大语言模型智能体的交互与上下文管理，使其在多轮、长程交互中持续维护用户意图与任务目标。
\section*{教育经历}

\noindent\textbf{香港科技大学（HKUST）}\hfill 2025 年 9 月至今\\
计算机科学与工程博士（在读）；导师：周晓方教授\par\vspace{3pt}
\noindent\textbf{哈尔滨工业大学}\hfill 2021 年 9 月 - 2025 年 6 月\\
计算机科学与技术学士，未来技术学院（原英才学院）\par\vspace{3pt}
\section*{研究方向}

大语言模型智能体的交互与上下文管理，重点包括：
\begin{itemize}[leftmargin=12pt,itemsep=1pt,topsep=2pt]
\item \textbf{用户交互}: 跟踪多轮对话中用户意图与偏好的变化。
\item \textbf{环境交互}: 在长程执行中持续维护用户意图与任务目标。
\end{itemize}
\section*{发表论文}

\begin{enumerate}[label={[\arabic*]},leftmargin=18pt,itemsep=1.2pt,topsep=1pt]
\item \textbf{Junle Chen}, Wei Chen, Zhengjun Huang, Zhoujin Tian, Yuxuan Liu, Kai Wang, Rui Chen, and Xiaofang Zhou. \href{https://arxiv.org/abs/2610.06496v1}{\textbf{You Changed Your Mind, The Model Didn't: Demystifying Intent in Multi-Turn Dialogue}}. arXiv:2610.06496, 2026.
\item \textbf{Junle Chen}, Wei Chen, Yehong Xu, Zhengjun Huang, Yuqian Wu, Zhoujin Tian, Kai Wang, Rui Chen, and Xiaofang Zhou. \href{https://arxiv.org/abs/2606.21169}{\textbf{Trip+: Benchmarking Agents in Personalized Interactive Travel Planning}}. \textbf{EMNLP 2026 (Main)}.
\item Wei Chen, \textbf{Junle Chen}, Yitong Yang, Zhaoyang Xu, Jiaxin Lin, Yuxuan Liang, Xiaofang Zhou, Kai Wang, and Rui Chen. \href{https://arxiv.org/abs/2610.06105}{\textbf{Flash-OPD: Fast On-Policy Distillation}}. arXiv:2610.06105, 2026.
\item Yuxuan Liu, Haoran Li, Yuhao Zhang, Jiahe Guo, Hongyu Luo, Wenbin Hu, Huihao Jing, Kawai Chung, \textbf{Junle Chen}, Changxuan Fan, Qing Zong, Lingyun Xie, and Yangqiu Song. \href{https://arxiv.org/abs/2610.04008}{\textbf{SkillScriptBench: Benchmarking Self-Evolution of Executable Agent Skill Packages Beyond Markdown}}. arXiv:2610.04008, 2026.
\item Zhengjun Huang, Wenxuan Liu, Zhoujin Tian, Wei Chen, \textbf{Junle Chen}, Yuqian Wu, Fangyuan Zhang, Qintian Guo, and Xiaofang Zhou. \href{https://arxiv.org/abs/2606.07402}{\textbf{M\textsuperscript{3}Exam: Benchmarking Multimodal Memory for Realistic User-Agent Interactions}}. arXiv:2606.07402, 2026.
\item Yuqian Wu, Zhijie Deng, Wei Chen, Junwei Li, Yutian Jiang, \textbf{Junle Chen}, Zhengjun Huang, Qingxiang Liu, Jing Tang, Jiaheng Wei, and Yuxuan Liang. \href{https://arxiv.org/abs/2606.04660}{\textbf{LifeSide: Benchmarking Agents as Lifelong Digital Companions}}. arXiv:2606.04660, 2026.
\item Yuqian Wu, Wei Chen, Zhengjun Huang, \textbf{Junle Chen}, Qingxiang Liu, Kai Wang, Xiaofang Zhou, and Yuxuan Liang. \href{https://arxiv.org/abs/2604.11628}{\textbf{Back to Basics: Let Conversational Agents Remember with Just Retrieval and Generation}}. \textbf{COLM 2026}.
\item Wei Chen, Yuqian Wu, \textbf{Junle Chen}, Xiaofang Zhou, and Yuxuan Liang. \href{https://arxiv.org/abs/2602.20677}{\textbf{UrbanFM: Scaling Urban Spatio-Temporal Foundation Models}}. \textbf{KDD 2026 Oral}.
\item Wei Chen, \textbf{Junle Chen}, Yuqian Wu, Yuxuan Liang, and Xiaofang Zhou. \href{https://arxiv.org/abs/2602.19113}{\textbf{Learning from Complexity: Exploring Dynamic Sample Pruning of Spatio-Temporal Training}}. arXiv:2602.19113, 2026.
\item Jian Chen\textsuperscript{\(\dagger\)}, \textbf{Junle Chen}\textsuperscript{\(\dagger\)}, Donghua Yang, Guangyu Sui, Jinbao Wang, and Lina Chen. \href{https://doi.org/10.1007/s11280-025-01368-6}{\textbf{Adaptive Indexing for Coverage Queries over Transition Trajectories}}. \textbf{World Wide Web, 2025}.
\item Jian Chen, Hong Gao, Yuhong Shi, \textbf{Junle Chen}, Donghua Yang, and Jianzhong Li. \href{https://doi.org/10.1109/TKDE.2024.3514323}{\textbf{Maximizing Influence Query over Indoor Trajectories}}. \textbf{IEEE Transactions on Knowledge and Data Engineering (TKDE), 2025}.
\item Jian Chen, \textbf{Junle Chen}, Hong Gao, Donghua Yang, and Kaiqi Zhang. \href{https://doi.org/10.1007/978-981-97-7235-3_32}{\textbf{Efficient Coverage Query over Transition Trajectories}}. \textbf{APWeb-WAIM, 2024}.
\end{enumerate}
{\footnotesize \textsuperscript{\(\dagger\)} 共同第一作者。}\par
\section*{获奖情况}

\begin{tabularx}{\linewidth}{@{}X X@{}}
\textbullet\ 香港科技大学研究生助学金（PGS） & \textbullet\ 香港科技大学 RedBird 博士奖学金\\
\textbullet\ 国家奖学金（2022-2023 学年） & \textbullet\ 人民一等奖学金（每学期获得）\\
\textbullet\ 美国大学生数学建模竞赛 Finalist 奖 & \textbullet\ 大学生数学竞赛一等奖\\
\textbullet\ 省级优秀学生 & \textbullet\ 黑龙江省优秀毕业生\\
\textbullet\ 哈尔滨工业大学优秀毕业生 & \\
\end{tabularx}
\section*{科研经历}

\textbf{南京大学 LAMDA 实验室}\hfill 2023 年 7 月\\
数据驱动知识的人工智能云端科研项目；项目经历
\end{document}
